% 7. Formularze i obsługa błędów
\section{Formularze i obsługa błędów}\label{sec:formularze}

\begin{instrukcja}[wymagane]
Zasady projektowania formularzy z Wykładu 3: etykiety przy polach, kontrolki dobrane do danych,
tylko potrzebne pytania, sensowne wartości domyślne, Zlokalizowane Komunikaty O Błędach (które pole,
dlaczego, jak poprawić), wpisane dane nie znikają po błędzie.
\end{instrukcja}

\begin{table}[H]\centering
\caption{Pola formularza, walidacja i komunikaty}\label{tab:formularz}
{\small
\begin{tabularx}{\linewidth}{@{}p{2.6cm}p{2.8cm}p{3.6cm}L@{}}
\naglowektabeli \nt{Pole} & \nt{Kontrolka} & \nt{Reguła walidacji} & \nt{Komunikat o błędzie} \\
\wpisz{tytuł} & \wpisz{pole tekstowe} & \wpisz{wymagane, 3-120 znaków} & \wpisz{„Wpisz tytuł książki (co najmniej 3 znaki).”} \\
\wpisz{cena} & \wpisz{pole liczbowe} & \wpisz{0-500 zł} & \wpisz{...} \\
\wpisz{stan} & \wpisz{przełączniki (radio)} & \wpisz{domyślnie „dobry”} & \wpisz{-} \\
\wpisz{...} & & & \\
\bottomrule
\end{tabularx}}
\end{table}

\zrzut{formularz_blad.png}{Formularz z komunikatem o błędzie przy polu}{0.7}
\wpisz{Jak zabezpieczyliście czynności nieodwracalne (cofnięcie albo potwierdzenie).}
